\documentclass[11pt]{article}

\usepackage[margin=25mm]{geometry}

\usepackage{amsmath,amssymb,fontspec,unicode-math}

\IfFontExistsTF{DejaVu Serif}{\setmainfont{DejaVu Serif}}{\setmainfont{Latin Modern Roman}}

\IfFontExistsTF{Latin Modern Math}{\setmathfont{Latin Modern Math}}{\setmathfont{DejaVu Math TeX Gyre}}

\usepackage{xurl,hyperref}

\setlength{\parindent}{0pt}\setlength{\parskip}{6pt}

\begin{document}

\title{Path-Enhanced Contrastive Learning for Recommendation}

\author{ELVIA preliminary AI analysis}\date{21 October 2026}\maketitle

Preliminary assessment: Major.\par

Independent scientific review remains pending. Author adoption has not been requested.

The printed temporal sampler fails its own domain/order guarantee. Its contrastive derivative identities are wrong, while the rotation interpretation requires a normalization qualification.

\section*{Review scope}

Full supplied technical report read alongside the corresponding original-paper claim passages and applicable author-release source. This pass independently checks the stated reasoning and bounds; historical executions were not rerun.

\section*{Problem 1: Temporal transition is undefined and permits past edges}

Source: Section 3.3, Eqs.(8)–(12), PDF p.5

Assessment: supported; impact: major

\subsection*{Source claim (paraphrase)}

The sampler uses softmax logits \(w(u,v)/[t(u,v)-t_{\rm last}]\) over the full neighborhood and claims the resulting paths maintain temporal consistency.

\subsection*{Analysis}

Use one item connected to users at timestamps \(2,1,3\), each with weight \(1\). Follow a valid center prefix whose last timestamp is \(2\). Candidate gaps are \((0,-1,1)\), giving logits \((1/0,-1,1)\), so the printed real-valued transition is undefined. The reverse edge alone produces the zero gap on an undirected graph. Even excluding it leaves positive probability \(1/(1+e^2)\) of selecting timestamp \(1\), violating nondecreasing time.

\subsection*{Finding and scope}

The printed path sampler is neither well-defined on all allowed edges nor sufficient for its temporal-order guarantee. An unstated implementation filter could change the operational behavior, but is not part of the stated formula.

\subsection*{Proposed correction}

Mask to edges with \(t(u,v)\geq t_{\rm last}\), use a positive offset \(\delta\) in the time-gap denominator, and specify termination or restart when no admissible edge exists. Every retained transition then preserves nondecreasing timestamps.

The described correction remains a proposal; author adoption or a corresponding erratum is not established.

\section*{Problem 2: Contrastive derivatives use the wrong normalizers}

Source: Section 3.5, Eqs.(15)–(18), PDF pp.6–7; objectives (7) and (14)

Assessment: supported; impact: minor

\subsection*{Source claim (paraphrase)}

The displayed gradients give attraction and repulsion the same partition function, summing both positive and denominator exponentials.

\subsection*{Analysis}

For \(\ell(x)=-\log[A(x)/B(x)]\), direct differentiation gives \(-A^{\prime}(x)/A(x)+B^{\prime}(x)/B(x)\), not a common \(A+B\) denominator. With a self-positive anchor, positive candidate \(1\), negative candidate \(0\) and temperature \(1\), take \(A=e^{x^2}+e^x,B=A+1\). At \(x=0\), the true derivative is \(-1/2+1/3=-1/6\); the printed shared normalizer gives zero. A negative-only denominator does not rescue the formula.

\subsection*{Finding and scope}

The displayed derivative identities and their quantitative weighting interpretation fail. The scalar objectives remain well-defined, and unavailable training code might differentiate those objectives correctly.

\subsection*{Proposed correction}

Differentiate each log-sum with its own partition function, including candidate/self dependencies and every occurrence in the full loss.

The described correction remains a proposal; author adoption or a corresponding erratum is not established.

\section*{Problem 3: Rotation interpretation depends on normalization}

Source: Section 3.4, Eq.(13), PDF p.6; Appendix A.5, p.23

Assessment: not\_issue; impact: minor

\subsection*{Source claim (paraphrase)}

The paper calls normalized complex-coordinate multiplication a rotation without specifying global or coordinatewise normalization.

\subsection*{Analysis}

Writing the update as \(T_e(w)=w\odot\overline e\), one has \(\|T_e(w)\|^2=\sum_k|w_k|^2|e_k|^2\). This preserves every distance exactly when every \(|e_k|=1\). Global normalization fixes their squared sum at \(1\), incompatible with that condition in multiple complex coordinates. For real/imaginary halves \(w=(1,0,0,0),e=(1/2,1/2,1/2,1/2)\), the norm contracts from \(1\) to \(1/\sqrt2\).

\subsection*{Finding and scope}

Coordinatewise unit phases validate the multiplicative rotation interpretation. The unspecified normalization leaves a conditional scope qualification; no author-code error or recommendation-performance failure is established.

\section*{Limits}

Preliminary AI source/argument assessment; independent human scientific adjudication remains pending.\par

No author-provided or retained verification program was executed. Published empirical results were not independently reproduced.\par

\section*{Source identity}

NeurIPS 2025 official proceedings PDF (exact retained local copy inspected)

\url{https://proceedings.neurips.cc/paper_files/paper/2025/file/7cb33a355bbe68305e02a080ad4eb86a-Paper-Conference.pdf}

Reviewed PDF SHA-256: \nolinkurl{49e563ce4ba17ed8501f7c999e3f158bc2717b2a12f9e0287eba57a33663c927}

The SHA-256 identifies the exact retained PDF inspected in this rapid source/argument check. Original supplied report bytes are retained locally. The original PDF retrieval timestamp is taken from matching retained metadata.

\end{document}
